% Nota de cubierta independiente: dirige al lector de la Parte I hacia la Parte II.
% NO modifica el PDF publicado de la Parte I; pensada para anteponer / adjuntar / usar como nota.
\documentclass[11pt]{article}
\usepackage[utf8]{inputenc}
\usepackage[T1]{fontenc}
\usepackage{lmodern}
\usepackage[spanish,es-noquoting]{babel}
\usepackage{amsmath,amssymb,graphicx,hyperref}
\usepackage{xurl}
\usepackage[table]{xcolor}
\definecolor{ggteal}{HTML}{1B6F8C}
\definecolor{ggband}{HTML}{DCE7EB}
\usepackage[margin=1.1in]{geometry}
\usepackage{microtype}
\usepackage[most]{tcolorbox}
\usepackage{parskip}
\pagestyle{empty}
\hypersetup{colorlinks=true,linkcolor=ggteal,urlcolor=ggteal}
\newcommand{\lean}[1]{\texttt{#1}}

\begin{document}

{\centering
{\large\bfseries Nota al lector de la Parte~I}\\[2pt]
{\color{ggteal}\rule{0.5\textwidth}{0.8pt}}\\[6pt]
{\itshape La secuela propuesta al final de este art\'iculo ya est\'a escrita.}\\[14pt]
}

La Parte~I (\emph{De signos aleatorios a emparejamientos}) cerr\'o con una propuesta
y una promesa: demostrar el \textbf{teorema de Heilmann--Lieb} no mediante la
maquinaria de familias entrelazadas, sino a trav\'es del \textbf{\'arbol de caminos
de Godsil}, y---\emph{``si este art\'iculo tiene una secuela, esta es su primera
frase.''} Esa secuela ya existe.

\textbf{Parte~II --- \emph{Desplegar un grafo en un \'arbol: una demostraci\'on
verificada por m\'aquina del teorema de Heilmann--Lieb en Lean~4}} lleva a cabo la
propuesta por completo. Construye el \'arbol de caminos, demuestra la divisibilidad
$\mu_G \mid \mu_{T(G,u)}$ (cruzando el \'unico muro de elaboraci\'on en el que la
ruta del \'arbol de caminos se hab\'ia detenido), establece la identidad del bosque
y demuestra la cota de Ramanujan a partir de una estimaci\'on de Gershgorin
ponderado. Ambas mitades son \lean{sorry}-libres.

\begin{tcolorbox}[colframe=ggteal,colback=ggband!40,arc=3pt,boxrule=0.9pt]
\centering
Para todo grafo finito $G$, el polinomio de emparejamientos $\mu_G$ tiene todas sus
ra\'ices reales; y si $2\le\Delta$ y todo v\'ertice tiene grado $\le\Delta$, toda
ra\'iz de $\mu_G$ est\'a en
\[
  \bigl[-2\sqrt{\Delta-1},\;2\sqrt{\Delta-1}\bigr].
\]
\end{tcolorbox}

Este es exactamente el punto final que la Parte~I cartografi\'o como trabajo futuro
(``La siguiente piedra''): el primer Heilmann--Lieb verificado por m\'aquina, que
completa el programa del \'arbol de caminos. Como en la Parte~I, la matem\'atica es
cl\'asica y no se reclama ning\'un teorema nuevo; la aportaci\'on es la
formalizaci\'on.

\bigskip
{\color{ggteal}\rule{\textwidth}{0.4pt}}\\[4pt]
\textbf{D\'onde encontrar la Parte~II.}
\begin{itemize}\setlength{\itemsep}{1pt}
\item Fuentes y ambos PDFs (ingl\'es + espa\~nol):
  \url{https://github.com/karlesmarin/godsil-gutman-lean}
\item Archivada en Zenodo: DOI \texttt{10.5281/zenodo.20561832}
  (\url{https://doi.org/10.5281/zenodo.20561832}).
\end{itemize}

\vfill
{\footnotesize\color{gray}
Carles Mar\'in, \texttt{karlesmarin@gmail.com}. Esta nota acompa\~na, y no altera, la
Parte~I archivada (DOI de Zenodo \texttt{10.5281/zenodo.20517350}).}

\end{document}
